% Generated draft: requires mathematical and source review.
\documentclass[11pt]{article}
\usepackage[margin=25mm]{geometry}
\usepackage{fontspec}
\setmainfont{Latin Modern Roman}
\usepackage{amsmath,amssymb,mathtools,unicode-math}
\setmathfont{Latin Modern Math}
\usepackage{enumitem,needspace,xurl,hyperref}
\setlength{\parindent}{0pt}
\setlength{\parskip}{5pt}
\setlength{\emergencystretch}{3em}
\newcommand{\N}{\mathbb{N}}
\newcommand{\Z}{\mathbb{Z}}
\newcommand{\Q}{\mathbb{Q}}
\newcommand{\R}{\mathbb{R}}
\newcommand{\C}{\mathbb{C}}
\newcommand{\F}{\mathbb{F}}
\newcommand{\A}{\mathbb{A}}
\newcommand{\PP}{\mathbb{P}}
\newcommand{\E}{\mathbb{E}}
\newcommand{\eps}{\varepsilon}
\newcommand{\dd}{\,\mathrm{d}}
\begin{document}
\section*{UnsolvedMath ID: 30006281}
\textbf{Source problem code:} catalogue entry
\subsection{Definitions and mathematical statement}
Write \(D^b(W)=D^b(\operatorname{Coh}W)\) for the bounded derived category of coherent sheaves on a smooth variety \(W\). An ordered semiorthogonal decomposition \(D^b(X)=\langle\mathcal A_1,\ldots,\mathcal A_m\rangle\) consists of full triangulated subcategories generating \(D^b(X)\), with \(\operatorname{Hom}(A_j,A_i[k])=0\) for \(j>i\), every \(k\in\Z\), and objects \(A_j\in\mathcal A_j\), \(A_i\in\mathcal A_i\). Mutations of adjacent admissible components give the categorical braid action referred to in the source.

Let \(X\) be a smooth Fano variety admitting extremal contractions \(f:X\to Y\) and \(g:X\to Z\), where \(Y\) and \(Z\) are smooth Fano varieties. Suppose that the semiorthogonal decompositions associated with these contractions are compatible with pullback and are related by a mutation braid \(\beta_{\mathrm{mut}}\), as intended in the source's Conjecture 3. That conjecture is not included in the supplied record, so neither the precise decompositions nor their pullback compatibility nor the choice of braid is specified here.

For a quantum parameter \(\tau\) and a small real B-field perturbation \(B\), put \(t=\tau+iB\) and denote quantum multiplication by the first Chern class by
\[
M_t:\operatorname{QH}^{*}(X,\C)_t\longrightarrow\operatorname{QH}^{*}(X,\C)_t,
\qquad a\longmapsto c_1(X)\star_t a.
\]
This notation presupposes a specialization of small quantum cohomology at which the product is defined; the source does not specify its analytic domain or parameter convention. Consider a path \(\gamma(s)=\tau(s)+iB(s)\), with \(\tau(s)\) in the ample cone of \(X\), approaching boundary walls associated with the contractions. The source proposes choosing the small B-field perturbation to avoid collisions of eigenvalues along the path.

Where \(M_t\) varies analytically and has \(d\) distinct eigenvalues, with \(d\) the dimension of its underlying complex vector space, its unordered spectrum lies in
\[
\operatorname{UConf}_d(\C)=\{\{z_1,\ldots,z_d\}\subset\C:z_i\ne z_j\text{ for }i\ne j\}.
\]
Tracking this spectrum along \(\gamma\) gives a spectral braid path. A loop of unordered spectra gives a braid in \(B_d=\pi_1(\operatorname{UConf}_d(\C))\), after a base configuration is chosen; a general open path gives a morphism in the corresponding fundamental groupoid. The existence of a suitable distinct-spectrum locus is not supplied as a hypothesis or established in the record.

\textbf{Conjecture 4, as represented by the supplied record.} The mutation braid relating the decompositions associated with \(f\) and \(g\) can be computed by tracking the eigenvalues of \(M_{\tau+iB}\) along the indicated path in the ample cone. After the required categorical and spectral identifications, the intended assertion is
\[
\beta_{\mathrm{mut}}=\beta_{\mathrm{quantum}}.
\]
As the path approaches the walls, the eigenvalues are expected to cluster into groups reflecting the corresponding birational contractions. The source does not specify the clustering rule, the path and endpoint identifications, or an identification of the categorical braid group with a braid group on eigenvalues or clusters. Consequently, this displayed equality records the intended conjectural comparison and is not yet a fully specified equality of elements of a fixed group. The base field of the varieties is also not explicit in the supplied statement, although the quantum cohomology has coefficients in \(\C\).
\subsection{Short English statement}
For a smooth Fano variety admitting extremal contractions to smooth Fano varieties \(Y\) and \(Z\), the conjecture proposes recovering the mutation braid between the associated derived-category decompositions from the motion of the eigenvalues of quantum multiplication by \(c_1(X)\). The quantum parameter moves toward walls of the ample cone with a small B-field perturbation intended to prevent collisions. The missing Conjecture 3 and unspecified path, clustering, and braid identifications prevent a fully precise formulation from this record alone.
\subsection{Sources}
\begin{itemize}
\item Supplied local UnsolvedMath registry record: permanent ID \(30006281\), exact source problem code \texttt{catalogue entry}, title \emph{Quantum-Monodromy Computation of Mutation Braids}. The record attributes the question to Conjecture 4 in an Oberwolfach workshop report, assigns the year 2025, and names no individual proposer.
\item \emph{Recent Trends in Algebraic Geometry} (2025), Oberwolfach Reports, pp.\ 1451--1516, DOI \texttt{10.4171/owr/2025/28}: \url{https://doi.org/10.4171/owr/2025/28}. This is also the supplied literature-source link labelled ``original OWR report.'' No more specific location within the report is supplied.
\item Literature assessment dated exactly \texttt{2026-08-22}: the supplied curation labels the problem open and describes the identification as explicitly conjectural in general, with example-level evidence. This is supplied dated curation, not an independently verified literature assessment. The background expressly requests independent verification of time-sensitive status and cited claims.
\item The registry labels its statement \texttt{corrected\_verified} and reports that the official report was checked, the question isolated from surrounding exposition, and notation repaired. These are supplied verification claims; this draft does not independently verify them.
\item The excerpt disclaims originality for the conjectures and points to references [4--7, 13] for related literature, [15] for numerous braid comparisons, [16] for weaving patterns controlling Fourier--Mukai kernels, and [11] for a result attributed to Mu\~noz. Bibliographic entries for these reference labels are not supplied, so their identities and claims remain unverified.
\end{itemize}
\end{document}
